
\documentclass{article}
\usepackage{colortbl}
\usepackage{makecell}
\usepackage{multirow}
\usepackage{supertabular}

\begin{document}

\newcounter{utterance}

\centering \large Interaction Transcript for game `eqbench', experiment `eq\_bench', episode 146 with reflexion{-}Qwen3.5{-}27B.
\vspace{24pt}

{ \footnotesize  \setcounter{utterance}{1}
\setlength{\tabcolsep}{0pt}
\begin{supertabular}{c@{$\;$}|p{.15\linewidth}@{}p{.15\linewidth}p{.15\linewidth}p{.15\linewidth}p{.15\linewidth}p{.15\linewidth}}
    \# & \multicolumn{2}{c}{Player} && \multicolumn{2}{c}{Game Master} \\
    \hline

    \theutterance \stepcounter{utterance}  
    & & & \multicolumn{4}{p{0.6\linewidth}}{
        \cellcolor[rgb]{0.9,0.9,0.9}{
            \makecell[{{p{\linewidth}}}]{
                \texttt{\tiny{[P1$\langle$GM]}}
                \texttt{You have attempted this same task before and did not succeed. These are the plans you wrote after those attempts.} \\
\texttt{Trial \#1: New plan: My previous attempt failed because I assigned a high score to "Betrayed" without sufficient textual evidence, as the dialogue focuses on rivalry and past grievances rather than a specific act of betrayal by Alejandro against Jackson. I also underestimated "Embarrassed," ignoring the public nature of the insults ("making a spectacle," "parading around") which would likely cause shame. Next time, I will strictly base scores on explicit dialogue cues: I will lower "Betrayed" to 0 or 1 since no betrayal is mentioned, and significantly raise "Embarrassed" to reflect the humiliation of being publicly shamed about his relationship and character. I will also re{-}evaluate "Guilt" based on the specific accusation of "stepping on others," though the dialogue suggests Jackson denies this, keeping that score low. Finally, I will ensure "Apathetic" remains 0 given the high emotional intensity of the exchange.} \\
\\ 
\texttt{Trial \#2: New plan: My previous attempt failed because I misidentified the primary emotional driver of the final exchange. While "Embarrassed" was a reasonable guess, the dialogue's climax ("parading around with another man's woman" vs. "couldn't keep her") is a direct attack on Jackson's masculinity and his failure to retain a partner, which aligns more strongly with **Guilt** (internalizing the failure) or **Betrayed** (if he views the woman's departure as a betrayal), but most critically, I likely undervalued the specific sting of the "couldn't keep her" line. The previous plan from Trial \#1 correctly identified the need to lower "Betrayed" and raise "Embarrassed," but the execution in the failed attempt still resulted in a loss, suggesting the scoring distribution was still off.} \\
\\ 
\texttt{Next time, I will re{-}evaluate the specific accusation "I ain't the one who couldn't keep her." This implies Jackson feels a deep sense of personal failure and shame regarding the loss of the woman, which is a form of **Guilt** or profound **Embarrassment**. However, looking at the options, the most intense reaction to being told you "couldn't keep" someone you loved is often a mix of shame and a feeling of being wronged. I will adjust my strategy to:} \\
\texttt{1.  **Re{-}evaluate "Guilt":** The line "couldn't keep her" directly attacks his capability and implies he is at fault for the loss. I will assign a higher score here (e.g., 7{-}8) as he is being forced to confront his failure.} \\
\texttt{2.  **Re{-}evaluate "Embarrassed":** The public nature of the insults ("spectacle," "parading") supports a high score, but the specific "couldn't keep her" line might push the emotion more toward internal shame/guilt than just public embarrassment. I will keep this high but perhaps slightly lower than Guilt if the prompt implies internal feeling.} \\
\texttt{3.  **Re{-}evaluate "Betrayed":** While the woman left, the dialogue doesn't explicitly state she betrayed him, but rather that he failed to keep her. I will keep this low (0{-}2) as the text focuses on his failure, not her treachery.} \\
\texttt{4.  **Re{-}evaluate "Apathetic":** Keep at 0 due to the high aggression.} \\
\\ 
\texttt{Specific Action: I will assign **Guilt: 8** (due to the "couldn't keep her" accusation), **Embarrassed: 6** (due to the public spectacle), **Betrayed: 1** (minimal evidence), and **Apathetic: 0**. This shifts the focus from the public shame to the personal failure highlighted in the final line.} \\
\\ 
\texttt{Trial \#3: New plan: My previous attempts failed because I misinterpreted the core emotional dynamic of the final exchange. The dialogue is a heated rivalry where Jackson attacks Alejandro's character ("conscience," "stepping on others") and Alejandro counters by attacking Jackson's masculinity and relationship status ("couldn't keep her"). The key mistake was assigning high scores to "Guilt" (7) and "Embarrassed" (8). Jackson is portrayed as aggressive and defensive ("I don't recall stepping on others," "I ain't the one..."), not as someone internalizing shame or admitting fault. Therefore, "Guilt" should be near zero because Jackson explicitly denies wrongdoing. "Embarrassed" is also likely low because he is actively fighting back rather than shrinking away. The most intense emotion for Jackson, given the accusation of "parading around with another man's woman," is **Betrayed** (feeling wronged by the woman's choice and Alejandro's possession of her) or perhaps a mix of anger (not an option) and resentment. However, looking at the specific options, the prompt likely expects the user to recognize that Jackson feels **Betrayed** by the situation (the woman leaving him for Alejandro) and potentially **Embarrassed** by the public nature of the "couldn't keep her" jab, but the *dominant* feeling in a "victim" narrative (as Alejandro calls him) is often a sense of being wronged.} \\
\\ 
\texttt{Wait, let's re{-}read the specific failure. The previous attempt gave Betrayed: 1, Embarrassed: 8, Guilt: 7. The outcome was LOSE.} \\
\texttt{The dialogue:} \\
\texttt{Jackson: "I ain't the one parading around with another man's woman." {-}> He feels Alejandro stole his woman. This is a classic **Betrayal** scenario (by the woman and Alejandro).} \\
\texttt{Alejandro: "And I ain't the one who couldn't keep her." {-}> This is a counter{-}attack.} \\
\texttt{Jackson's reaction to "couldn't keep her" is likely **Anger** (not an option) or **Embarrassment** (shame at failure).} \\
\texttt{But the prompt asks for "Betrayed, Embarrassed, Guilt, Apathetic".} \\
\texttt{If Jackson feels he was "stolen from," **Betrayed** should be the highest score. My previous plan gave it a 1. That was the error. I focused too much on the "couldn't keep her" line as a source of guilt/shame, but Jackson's stance is "You stole her," which is betrayal.} \\
\texttt{Also, "Guilt" is definitely low because Jackson is defensive ("I don't recall...").} \\
\texttt{"Embarrassed" is possible due to the public nature, but "Betrayed" is the root cause of the conflict.} \\
\\ 
\texttt{Corrected Strategy:} \\
\texttt{1.  **Betrayed**: This must} \\
\\ 
\\ 
\\ 
\texttt{Your task is to predict the likely emotional responses of a character in this dialogue:} \\
\\ 
\texttt{Jackson: Sanchez. Never thought I'd see you here, thought you'd be too busy basking in your glory.} \\
\texttt{Alejandro: Riley, always the charmer. I could say the same about you. How's the knee?} \\
\texttt{Jackson: Better than your conscience, I reckon.} \\
\texttt{Alejandro: Ah, the ever{-}righteous Riley. You think you're the only one who's had to fight for everything?} \\
\texttt{Jackson: I don't recall stepping on others to get to the top.} \\
\texttt{Alejandro: And I don't recall asking for your opinion.} \\
\texttt{Jackson: You never did. Too busy making a spectacle of yourself.} \\
\texttt{Alejandro: And you're too busy playing the victim. } \\
\texttt{Jackson: I ain't the one parading around with another man's woman.} \\
\texttt{Alejandro: And I ain't the one who couldn't keep her.} \\
\texttt{{[}End dialogue{]}} \\
\\ 
\texttt{At the end of this dialogue, Jackson would feel...} \\
\texttt{Betrayed} \\
\texttt{Embarrassed} \\
\texttt{Guilt} \\
\texttt{Apathetic} \\
\\ 
\texttt{Give each of these possible emotions a score from 0{-}10 for the relative intensity that they are likely to be feeling each.} \\
\\ 
\texttt{You must output in the following format, including headings (of course, you should give your own scores), with no additional commentary:} \\
\\ 
\texttt{Betrayed: <score>} \\
\texttt{Embarrassed: <score>} \\
\texttt{Guilt: <score>} \\
\texttt{Apathetic: <score>} \\
\\ 
\\ 
\texttt{{[}End of answer{]}} \\
\\ 
\texttt{Remember: zero is a valid score, meaning they are likely not feeling that emotion. You must score at least one emotion > 0.} \\
\\ 
\texttt{Your answer:} \\
            }
        }
    }
     \\ \\

    \theutterance \stepcounter{utterance}  
    & \multicolumn{4}{p{0.6\linewidth}}{
        \cellcolor[rgb]{0.561,0.737,0.561}{
            \makecell[{{p{\linewidth}}}]{
                \texttt{\tiny{[P1$\rangle$GM]}}
                \texttt{Betrayed: 9} \\
\texttt{Embarrassed: 7} \\
\texttt{Guilt: 1} \\
\texttt{Apathetic: 0} \\
            }
        }
    }
    & & \\ \\

    \theutterance \stepcounter{utterance}  
    & & & \multicolumn{2}{p{0.3\linewidth}}{
        \cellcolor[rgb]{0.9,0.9,0.9}{
            \makecell[{{p{\linewidth}}}]{
                \texttt{\tiny{[GM$|$GM]}}
                \texttt{Betrayed: 9} \\
\texttt{Embarrassed: 7} \\
\texttt{Guilt: 1} \\
\texttt{Apathetic: 0} \\
            }
        }
    }
    & & \\ \\

    \theutterance \stepcounter{utterance}  
    & & & \multicolumn{2}{p{0.3\linewidth}}{
        \cellcolor[rgb]{0.9,0.9,0.9}{
            \makecell[{{p{\linewidth}}}]{
                \texttt{\tiny{[GM$|$GM]}}
                \texttt{Betrayed: 8} \\
\texttt{Embarrassed: 6} \\
\texttt{Guilt: 0} \\
\texttt{Apathetic: 0} \\
            }
        }
    }
    & & \\ \\

    \theutterance \stepcounter{utterance}  
    & & & \multicolumn{2}{p{0.3\linewidth}}{
        \cellcolor[rgb]{0.9,0.9,0.9}{
            \makecell[{{p{\linewidth}}}]{
                \texttt{\tiny{[GM$|$GM]}}
                \texttt{game\_result = LOSE} \\
            }
        }
    }
    & & \\ \\

    \theutterance \stepcounter{utterance}  
    & & \multicolumn{4}{p{0.6\linewidth}}{
        \cellcolor[rgb]{0.95,0.95,0.95}{
            \makecell[{{p{\linewidth}}}]{
                \texttt{\tiny{[P1$\leftrightarrow$P]}}
                \texttt{Post{-}episode Reflexion (not sent to GM):} \\
\texttt{New plan: My previous attempts failed because I misidentified the primary emotional driver of Jackson's final reaction. I incorrectly prioritized "Guilt" and "Embarrassed" based on the "couldn't keep her" line, interpreting it as Jackson internalizing failure. However, Jackson's dialogue ("I ain't the one parading around with another man's woman") explicitly frames the situation as Alejandro stealing his partner, which is a classic **Betrayal** scenario. Jackson is defensive and accusatory, not remorseful, so "Guilt" should be near zero. While the public nature of the argument suggests some "Embarrassment," the core emotion driving his aggression is the feeling of being wronged by the woman and Alejandro.} \\
\\ 
\texttt{Next time, I will strictly align the scores with Jackson's accusatory stance:} \\
\texttt{1.  **Betrayed**: Assign the highest score (9{-}10) because the entire conflict stems from the accusation that Alejandro is with "another man's woman," implying a betrayal of Jackson's relationship.} \\
\texttt{2.  **Embarrassed**: Assign a moderate score (4{-}6) to acknowledge the sting of the "couldn't keep her" insult in a public setting, but keep it secondary to the feeling of betrayal.} \\
\texttt{3.  **Guilt**: Assign a very low score (0{-}1) because Jackson explicitly denies wrongdoing ("I don't recall stepping on others") and attacks Alejandro instead.} \\
\texttt{4.  **Apathetic**: Assign 0, as the dialogue is highly emotional and aggressive.} \\
\\ 
\texttt{Specific Action: I will output **Betrayed: 9**, **Embarrassed: 5**, **Guilt: 0**, and **Apathetic: 0**. This corrects the error of undervaluing the betrayal aspect and overvaluing the guilt aspect.} \\
            }
        }
    }
    & \\ \\

\end{supertabular}
}

\end{document}
